\textit{Statistical Mechanics.} Our model draws upon the statistical
mechanics of interacting spins: the energy function is the Ising model~\citep{ising1925}, the opinion update is its Glauber dynamics~\citep{glauber1963}, and the order parameter separating polarization from consensus is the Edwards--Anderson parameter of
spin glasses~\citep{edwards1975,sherrington1975}. The idea that such energy-based systems perform collective computation originates with Hopfield networks~\citep{hopfield1982neural,amit1985storing}.


\textit{Applications of Statistical Mechanics.} The Ising model characterizes the phases of spin glasses, where competing couplings accommodate many distinct low-energy states~\citep{parisi1979,mezard1987}, and of flocks, where local alignment among moving agents produces collective order~\citep{vicsek1995,toner1995}. Modern work has shown that the same theory can provide descriptive models of neural~\citep{schneidman2006}, protein-sequence~\citep{weigt2009,marks2011}, and collective-behavior~\citep{bialek2012statistical} data, and also emerges naturally from the theory of binary choice under social influence~\citep{brock2001,castellano2009}.


\textit{Social Simulations.} A growing literature uses LLMs to simulate people and societies, including experimental
subjects~\citep{argyle2023,horton2023,park2024}, populations within virtual communities~\citep{park2023,yang2024oasis,piao2025}, and systems exhibiting collective phenomena such as opinion dynamics~\citep{chuang2024,ashery2025,demarzo2024}.
These works demonstrate that LLM communities can represent complex behaviors; we offer a quantitative, \textit{predictive} theory of how they evolve.

\textit{Multi-agent Systems.} Modern multi-agent systems leverage interaction among LLM agents to improve reasoning, often
through debate or role-conditioned discussion~\citep{du2023improvingfactualityreasoninglanguage,liang-etal-2024-encouraging,chen2023, Swanson2024.11.11.623004}. Though early work manually specified interactions, these systems are increasingly designed automatically by searching over prompts, workflows, and graphs~\citep{khattab2023,zhuge2024,hu2025automateddesignagenticsystems, nielsen2026learningorchestrateagentsnatural}. Common multi-agent interaction patterns often rely on sub-task partitioning~\citep{yang2025agentnetdecentralizedevolutionarycoordination, zhang2025aflowautomatingagenticworkflow}, aggregation over individual agent opinions~\citep{chen2024reconcile, wang2024mixtureofagentsenhanceslargelanguage}, or open-ended deliberation~\citep{du2023improvingfactualityreasoninglanguage}. Moreover, systems that learn multi-agent interaction topologies often require expensive roll-outs, and the resulting systems can be brittle ~\citep{cemri2025,wang2024bounds}. 


